% Glossary
\chapter*{Glossary}
\addcontentsline{toc}{chapter}{Glossary}

\begin{description}

\item[Aliquot sequence] A sequence formed by iterating the sum-of-proper-divisors function: $s(n) = \sigma(n) - n$.

\item[American Mathematical Society (AMS)] The primary professional organization for mathematicians in the United States; publishes the \textit{Mathematics of Computation} and many reference texts cited in this book.

\item[Amicable pair] A pair of distinct integers $(a, b)$ such that the sum of the proper divisors of $a$ equals $b$ and the sum of the proper divisors of $b$ equals $a$.

\item[Bernoulli number] A sequence of rational numbers $B_n$ appearing in the Taylor series of $x/(e^x - 1)$ and in Euler's formula for $\zeta(2n)$.

\item[Chi ($\chi$)] The Euler characteristic, a topological invariant given by $V - E + F$ for polyhedra or $2 - 2g$ for surfaces of genus $g$.

\item[Chinese Postman Problem] The problem of finding the shortest closed path that visits every edge of a graph at least once; a generalization of the Eulerian circuit problem.

\item[Congruence] A relation $a \equiv b \pmod{n}$ meaning $n$ divides $a - b$; foundational to modular arithmetic and Euler's theorem.

\item[Cyclotomic polynomial] The polynomial $\Phi_n(x)$ whose roots are the primitive $n$-th roots of unity; its degree is $\varphi(n)$.

\item[Digamma function] The logarithmic derivative of the gamma function: $\psi(z) = \Gamma'(z)/\Gamma(z)$.

\item[Dirichlet distribution] A multivariate continuous probability distribution parameterized by a vector of positive reals; the digamma function appears in its normalization.

\item[Euler characteristic] The quantity $\chi = V - E + F$ for a polyhedral surface; for a sphere $\chi = 2$, for a torus $\chi = 0$.

\item[Euler numbers (secant)] The even-indexed coefficients in the Taylor series of $\sec x = \sum E_{2n} x^{2n}/(2n)!$: $E_0=1, E_2=1, E_4=5, E_6=61, E_8=1385, \ldots$

\item[Euler product] The identity $\zeta(s) = \prod_p (1 - p^{-s})^{-1}$ connecting the Riemann zeta function to the prime numbers.

\item[Euler's constant ($\gamma$)] Also called the Euler--Mascheroni constant, $\gamma = \lim_{n \to \infty} (H_n - \log n) \approx 0.5772156649$.

\item[Euler's formula] The identity $e^{ix} = \cos x + i \sin x$ linking the complex exponential to trigonometric functions.

\item[Euler's reflection formula] The identity $\Gamma(z)\Gamma(1-z) = \pi/\sin(\pi z)$ relating the gamma function at $z$ and $1-z$.

\item[Euler's theorem] The statement that if $\gcd(a,n) = 1$, then $a^{\varphi(n)} \equiv 1 \pmod{n}$.

\item[Euler zigzag numbers] Also called Andre permutations or up/down numbers. $A_n$ counts the alternating permutations of $\{1, \ldots, n\}$. Generating function: $\tan x + \sec x = \sum A_n x^n/n!$.

\item[Eulerian circuit] A closed walk in a graph that traverses every edge exactly once; exists iff every vertex has even degree.

\item[Eulerian path] A walk in a graph that traverses every edge exactly once; exists iff exactly 0 or 2 vertices have odd degree.

\item[Entringer triangle] A triangular array where row $n$ encodes the number of alternating permutations of $\{1, \ldots, n+1\}$ by their first element. Constructed by cumulative sums alternating direction (Seidel's method).

\item[Fermat's little theorem] The special case of Euler's theorem for prime modulus: if $p$ is prime and $\gcd(a,p) = 1$, then $a^{p-1} \equiv 1 \pmod{p}$.

\item[Fleury's algorithm] A greedy algorithm for finding an Eulerian path that avoids bridges unless forced; runs in $O(E^2)$ time.

\item[Fundamental Theorem of Arithmetic] Every integer greater than 1 has a unique factorization into prime numbers, up to ordering.

\item[Gaussian integral] The integral $\int_{-\infty}^{\infty} e^{-x^2} dx = \sqrt{\pi}$, used to prove $\Gamma(1/2) = \sqrt{\pi}$.

\item[Genus] The number of ``holes'' in a surface; a sphere has genus 0, a torus has genus 1.

\item[Graph] A mathematical structure consisting of vertices (nodes) connected by edges; used to model pairwise relationships.

\item[Hierholzer's algorithm] An efficient algorithm for finding an Eulerian circuit by decomposing the graph into cycles and splicing them together; runs in $O(E)$ time.

\item[Hamiltonian path] A path in a graph that visits every vertex exactly once; finding one is NP-complete.

\item[Harmonic number] The $n$-th partial sum of the harmonic series: $H_n = \sum_{k=1}^n 1/k$.

\item[Harmonic series] The divergent series $\sum_{n=1}^{\infty} 1/n = 1 + 1/2 + 1/3 + \cdots$.

\item[Identity (Euler's)] The equation $e^{i\pi} + 1 = 0$, connecting the five fundamental constants $0, 1, e, i, \pi$.

\item[Infinite product] An expression of the form $\prod_{n=1}^{\infty} a_n$ as the limit of partial products; Euler used them to represent trigonometric and zeta functions.

\item[Mersenne prime] A prime number of the form $2^p - 1$ where $p$ is also prime.

\item[Modular arithmetic] Arithmetic on equivalence classes modulo $n$, denoted $\mathbb{Z}/n\mathbb{Z}$; operations are performed on remainders.

\item[Perfect number] A positive integer equal to the sum of its proper divisors; equivalently, $\sigma(n) = 2n$.

\item[Platonic solid] One of the five convex regular polyhedra: tetrahedron, cube, octahedron, dodecahedron, icosahedron.

\item[Polyhedron] A three-dimensional solid bounded by flat polygonal faces, straight edges, and vertices.

\item[Primitive root] An integer $g$ modulo $p$ such that every nonzero residue modulo $p$ is a power of $g$; there are $\varphi(p-1)$ primitive roots modulo a prime $p$.

\item[Proper divisor] A divisor of $n$ excluding $n$ itself.

\item[Riemann zeta function] The function $\zeta(s) = \sum_{n=1}^{\infty} n^{-s}$ for $\Re(s) > 1$, extended to the complex plane by analytic continuation.

\item[Sociable number] A number belonging to a cycle of length greater than 2 in the aliquot sequence iteration.

\item[Sieve of Eratosthenes] An ancient algorithm for finding all primes up to $n$ by iteratively marking multiples of each discovered prime.

\item[Stirling's approximation] An asymptotic formula for the factorial: $n! \sim \sqrt{2\pi n}\,(n/e)^n$.

\item[Taylor series] An infinite series representation of a function as $\sum_{n=0}^{\infty} f^{(n)}(a)(x-a)^n/n!$; Euler used Taylor series extensively.

\item[Totient function ($\varphi$)] Euler's totient function counting the positive integers up to $n$ that are coprime to $n$.

\item[Vertex] A point or node in a graph or polyhedron; in graph theory, vertices are connected by edges.

\item[Wallis product] An infinite product for $\pi/2$: $\frac{\pi}{2} = \prod_{n=1}^{\infty} \frac{4n^2}{4n^2 - 1}$.

\item[Zeta function value at even integers] Euler's formula $\zeta(2n) = \frac{(-1)^{n+1} B_{2n} (2\pi)^{2n}}{2(2n)!}$ expressing even zeta values in terms of Bernoulli numbers and powers of $\pi$.

\item[Arbitrary-precision arithmetic] Computation using exact integer arithmetic scaled to a fixed number of decimal places, avoiding floating-point rounding errors.

\item[AGM (Arithmetic-Geometric Mean)] The common limit of the sequences $a_{n+1} = (a_n+b_n)/2$, $b_{n+1} = \sqrt{a_nb_n}$, used in the Brent-Salamin algorithm for computing $\pi$ quadratically.

\item[Binet formula] An integral representation of $\gamma$ involving an exponential kernel, enabling accelerated convergence.

\item[Chudnovsky algorithm] The fastest known algorithm for computing $\pi$, converging 14 digits per term, based on Ramanujan's series.

\item[Euler-Maclaurin summation] A technique for accelerating the convergence of series using correction terms derived from Bernoulli numbers.

\item[Machin's formula] The identity $\pi/4 = 4\arctan(1/5) - \arctan(1/239)$ used by Machin (1706) to compute 100 digits of $\pi$.

\item[Newton's method for sqrt] The iteration $x_{n+1} = (x_n + S/x_n)/2$ for computing $\sqrt{S}$, converging quadratically.

\end{description}

\cleardoublepage
